\section{Boundary classes and the character quotient}
\label{sec:descent}

We now apply Theorem~\ref{thm:diagram-density} to the surface underlying
the curve.  Flags encode the boundary classes by closed rank bounds.
We then make the handle determinants one and select the exact classes
on the character quotient.  This order is important: semisimplifying
a representation can decrease the Jordan blocks of a boundary
matrix.

\subsection{Flags for arbitrary Jordan classes}

Let $\overline X$ be the smooth compactification of the curve, of
genus $g$, and replace each missing point by a boundary circle.
For a prescribed boundary class, write $\nu_\lambda$ for its Jordan
partition at the eigenvalue $\lambda$, and write
\[
 c_{\lambda,1}\ge c_{\lambda,2}\ge\cdots>0
\]
for the column lengths of this partition.  Its generalized
$\lambda$-eigenspace has dimension $n_\lambda=|\nu_\lambda|$.
Order the eigenvalues once and for all.  On the boundary put a flag
whose graded pieces, in that order, have dimensions
$c_{\lambda,1},c_{\lambda,2},\ldots$ and scalar monodromy $\lambda$.
Realize each piece by a punctured disk cap, choosing the scalar
according to the meridian orientation.

Every matrix in the prescribed class admits this flag: on its
generalized $\lambda$-eigenspace use the successive kernels of powers
of $A-\lambda I$.  Conversely, any matrix with the prescribed flag
has the required characteristic polynomial and satisfies
\begin{equation}\label{eq:jordan-rank-bounds}
 \rk\bigl((A-\lambda I)^k\bigr)
 \le r-\sum_{j=1}^k c_{\lambda,j}
 \qquad(\lambda\text{ prescribed},\ k\ge1),
\end{equation}
where missing columns are zero.  To see this, restrict the flag to
the generalized $\lambda$-eigenspace.  The nilpotent part lowers its
flag by one step, so the first $k$ steps lie in its $k$th kernel.
On the other generalized eigenspaces, $A-\lambda I$ is invertible.
These observations give~\eqref{eq:jordan-rank-bounds}.
Restriction to the generalized eigenspaces is valid because their
polynomial projectors preserve every invariant flag.  This argument
takes place over $\C$ and requires no unit condition on differences
between the prescribed boundary eigenvalues.

Use the compact surface as a rank-$r$ parent facet and attach the
caps just described.  In the projective case no caps are needed.
Let $Z$ be the reduced Zariski closure, in the framed $\GL_r$
representation variety, of the main-surface representations obtained
from this diagram.  Theorem~\ref{thm:diagram-density} gives a dense
set of actual $\GL_r(R)$ matrix points in $Z$.  The preceding flag
argument shows that $Z$ contains every representation with the
prescribed exact classes and lies within the closed conditions
consisting of their characteristic polynomials and
\eqref{eq:jordan-rank-bounds}.  In particular all its peripheral
determinants are one.

The set $Z$ is invariant under simultaneous conjugation.  It is also
invariant under twisting by characters of $\pi_1(\overline X)$:
such a twist independently multiplies each of the $2g$ handle
matrices by a scalar and leaves all peripheral and cap data
unchanged.  These two invariances also hold for its closure.

\subsection{A single field makes all determinants one}

Put $S=Z\cap\Hom(\pi_1(X),\SL_r)$, taking reduced varieties, and
let $T=(\mathbb G_m)^{2g}$.  Write $\delta:Z\to T$ for the handle
determinants.  Scalar twisting defines
\[
 f:S\times T\longrightarrow Z,\qquad (s,t)\longmapsto t\cdot s.
\]
It is the base change of the power map $[r]:T\to T$.  More
explicitly,
\begin{equation}\label{eq:determinant-pullback}
 \begin{split}
 S\times T&\ \xrightarrow{\ \sim\ }\
       \{(z,t)\in Z\times T:\delta(z)=t^r\},\\
 (s,t)&\longmapsto(t\cdot s,t),
 \end{split}
\end{equation}
with inverse $(z,t)\mapsto(t^{-1}\cdot z,t)$.
Twist invariance of $Z$ ensures that this inverse takes values in
$S\times T$.  Consequently $f$ is finite, \'etale and surjective,
and in particular open.

Let $L/F$ be the fixed finite extension of
Lemma~\ref{lem:uniform-field}, containing all $r$th roots of all
units of $R$.  If $z$ is one of the integral matrix points dense in
$Z$, every coordinate of $\delta(z)$ is a unit of $R$.
Every preimage of $z$ under $f$ therefore has $t\in
(\cO_L^\times)^{2g}$ and $s=t^{-1}\cdot z$ an actual
$\SL_r(\cO_L)$ representation.  Since $f$ is open, the full
preimage of this dense set is dense in $S\times T$.  Projection to
$S$ shows that actual $\SL_r(\cO_L)$ representations are dense
in $S$.  The field $L$ is chosen before the points and is the same
for all components.  For $g=0$, the torus $T$ is trivial and this
argument is the identity map.

\subsection{Selecting the exact classes after semisimplification}

Let $\mathcal R=\Hom(\pi_1(X),\SL_r)$ over $\C$ and
let $q:\mathcal R\to M_r(X)_\C$ be its affine character quotient.
Let $C\subset\mathcal R$ be the invariant closed locus defined by
the prescribed characteristic polynomials and all the bounds
\eqref{eq:jordan-rank-bounds}.  Only $1\le k\le r$ is needed.
Let $C_{\mathrm{bad}}\subset C$ be the union of the loci where at
least one rank is strictly smaller than the rank in the prescribed
class.  This is a finite union of invariant closed subsets, obtained
by lowering one of the positive rank bounds by one.

Lemma~\ref{lem:exact-stratum} identifies the exact character locus as
\begin{equation}\label{eq:exact-character-open}
             Y=q(C)\setminus q(C_{\mathrm{bad}}).
\end{equation}
In particular it is open in $q(C)$, including when semisimplification
can shrink some boundary blocks.

We have $S\subset C$, and $S$ contains every representation with
the exact prescribed boundary classes.  Hence $q(S)$ is a closed
subset of $q(C)$ containing $Y$.  The dense integral representations
in $S$ give a dense subset $A\subset q(S)$ of ambient integral
character points: evaluating integral conjugation invariants on
their entries gives points of the specified model over $\cO_L$.
No integral matrix realization of their semisimplifications is
needed.  By~\eqref{eq:exact-character-open}, $Y$ is open in $q(S)$,
so $A\cap Y$ is dense in $Y$.

This proves Theorem~\ref{thm:main} for exact classes.  It proves
density in the entire reduced locus, hence in every irreducible
component, without an irreducibility assumption on $Z$, $S$ or $Y$.
It also retains the genus-zero and nonregular boundary cases.
The selection concerns only the generic-fiber character; it imposes
no avoidance of $q(C_{\mathrm{bad}})$ at finite primes.

Finally, fixed characteristic polynomials allow only finitely many
Jordan partitions.  Their exact strata cover the fixed-polynomial
locus.  The same $F$ and $L$ work for all of them, because these
fields depend on the rank and the boundary eigenvalues, not on the
partitions.  Taking the finite union of the dense sets proves the
fixed-polynomial assertion as well.  Empty strata are harmless.

For completeness, if $X$ is nonprojective and no boundary conditions
are imposed, its fundamental group is free of rank $2g+s-1$.
The group $\SL_r(\Z)$ is Zariski dense in $\SL_r(\C)$: its
closure contains every elementary one-parameter unipotent subgroup,
and these generate $\SL_r(\C)$.  Thus the integral matrix tuples
are dense in the free-group representation variety, and their
characters are dense in its quotient, already over $\Z$.
The cases of a free group of rank zero and of $r=1$ give a point.
Together with the projective case proved above, this establishes the
unrestricted assertion of Theorem~\ref{thm:main}.

\begin{remark}[The rank-two comparison]\label{rem:rank-two-comparison}
The ambient-integral exact-class statement in rank two also follows
from Coccia--Litt's surface results
\cite[Theorems~5.0.4 and~6.1.5]{CocciaLitt}.
Write $\varepsilon I$ for the product of the prescribed central
boundary matrices and remove those punctures from the presentation.
If another puncture remains, multiply its matrix by $\varepsilon$.
The relation becomes the ordinary surface relation, and the corresponding
trace changes only by a sign. These tuple maps have integral coefficients
and commute with conjugation; their pullbacks therefore induce maps on
integral invariant rings and preserve ambient integral points.

For the remaining noncentral boundary classes, the trace determines
the class unless it is $2$ or $-2$. In those cases exclude the closed
quotient locus of the corresponding scalar matrix. By
Lemma~\ref{lem:exact-stratum}, exactness is thus an open condition in
the trace locus, to which their ordinary surface theorem applies.
If every prescribed boundary matrix was central, the remaining
relation is
\[
                 \prod_{j=1}^g[A_j,B_j]=\varepsilon I.
\]
For $\varepsilon=1$ use their ordinary theorem; for
$\varepsilon=-1$ and $g\geq1$ use their twisted theorem. In genus
zero these two cases are respectively a point and the empty set.
This deduction imposes no exact Jordan condition at finite primes.
\end{remark}
